\documentclass{article}
\usepackage{geometry}
\geometry{a4paper, margin=1in}
\usepackage{hyperref}
\hypersetup{colorlinks=true, urlcolor=blue, citecolor=blue}

\title{Proposal C: Self-Supervised Cell Tracking via\\Temporal Cycle Consistency}
\author{Candidate Name}
\date{Suggested Supervisor: Prof.\ Dr.\ Martin Weigert}

\begin{document}
\maketitle

\section*{Motivation}

Cell tracking in time-lapse microscopy is bottlenecked by annotation cost. Current state-of-the-art methods---including Trackastra~\cite{trackastra}---are trained on manually annotated cell tracks: human experts label which cell in frame~$t$ corresponds to which cell in frame~$t{+}1$, including divisions and disappearances. For the Cell Tracking Challenge datasets, these annotations took months of expert effort per sequence. For dense 3D embryo data (e.g., \textit{Tribolium} with thousands of cells per frame across hundreds of timepoints), comprehensive track annotation is effectively infeasible.

What if the tracking model could learn \textbf{without any annotated tracks}?

The key insight is that cell identity is temporally consistent: if cell $A$ in frame~$t$ is matched to cell $B$ in frame~$t{+}1$, and $B$ is matched to cell $C$ in frame~$t{+}2$, then cycling back from $t{+}2$ to $t$ should recover $A$. This \textbf{cycle consistency} provides a self-supervised training signal. Similar ideas have transformed optical flow estimation (CycleFlow), point tracking (PIPs, CoTracker), and video correspondence learning---but have never been applied to cell tracking.

\section*{Approach}

\begin{enumerate}
	\item \textbf{Cell embedding network.} Given a segmented cell and its local image context, produce a feature embedding. Architecture follows Trackastra's design: extract image crops per cell, encode with a CNN or ViT backbone, output a fixed-dimensional embedding.

	\item \textbf{Soft association via embedding similarity.} For consecutive frames, compute a soft association matrix between all cell pairs using cosine similarity of their embeddings. This produces a differentiable ``tracking'' without hard assignments.

	\item \textbf{Cycle-consistency loss.} Track forward ($t \to t{+}1 \to t{+}2$) by composing soft association matrices, then track backward ($t{+}2 \to t{+}1 \to t$). The cycle-consistency loss penalizes deviation from the identity mapping. Cells that divide or disappear naturally violate cycle consistency, producing a learned signal for these events.

	\item \textbf{Auxiliary SSL objectives.} Combine cycle consistency with:
	      \begin{itemize}
		      \item \textbf{Temporal smoothness}: embeddings of the same cell should change slowly across time.
		      \item \textbf{Spatial distinctiveness}: embeddings of different cells in the same frame should be distinguishable.
		      \item \textbf{TAP signal}: optionally incorporate time-arrow prediction~\cite{tap} as a complementary pretext task on the shared encoder.
	      \end{itemize}

	\item \textbf{Evaluation.} Train on \textbf{unlabeled} Cell Tracking Challenge sequences and SLICE-1 \textit{Tribolium} data~\cite{slice1}. At test time, convert soft associations to hard assignments (Hungarian algorithm or greedy matching). Compare tracking quality (TRA, DET, AOGM metrics) against:
	      \begin{itemize}
		      \item Trackastra (fully supervised)
		      \item Simple overlap/nearest-neighbor baselines (no learning)
		      \item The self-supervised model fine-tuned on a small fraction (1\%, 5\%, 10\%) of annotated tracks
	      \end{itemize}
\end{enumerate}

\section*{Why This Works as a One-Semester Project}

\begin{itemize}
	\item \textbf{Trackastra provides the architectural template}: the cell crop extraction, association transformer, and evaluation pipeline are reusable. The change is in the training objective, not the architecture.
	\item \textbf{Cycle consistency is well-understood}: proven in optical flow, point tracking, and video correspondence. The adaptation to cell tracking is conceptually straightforward.
	\item \textbf{Strong baselines exist}: Trackastra (supervised) and simple overlap tracking (no learning) bracket the performance range. Any self-supervised result between these two is informative.
	\item \textbf{Incremental milestones}: (1)~reproduce Trackastra evaluation, (2)~implement cycle-consistency training on a simple 2D dataset, (3)~scale to embryo data, (4)~add auxiliary losses.
\end{itemize}

\section*{Expected Contribution}

The first \textbf{self-supervised cell tracking method}---learning to associate cells across time without any annotated tracks. Even if the method does not match fully supervised Trackastra, the scientifically interesting regime is \textbf{low-annotation fine-tuning}: pretrain self-supervised, then fine-tune on 1--10\% of annotated tracks. If this matches or exceeds training from scratch on 100\% of tracks, it demonstrates that SSL can dramatically reduce the annotation burden for cell tracking in embryogenesis.

\begin{thebibliography}{9}
	\bibitem{trackastra} B.~Gallusser and M.~Weigert. ``Trackastra: Transformer-based cell tracking for live-cell microscopy.'' \textit{ECCV}, 2024.
	\bibitem{tap} B.~Gallusser, M.~Stieber, and M.~Weigert. ``Self-supervised dense representation learning for live-cell microscopy with time arrow prediction.'' \textit{MICCAI}, 2023.
	\bibitem{slice1} F.~Strobl et al. ``SLICE-1: A comprehensive light-sheet imaging dataset of \textit{Tribolium castaneum} embryogenesis.'' \textit{Scientific Data}, Dec 2025.
\end{thebibliography}

\end{document}
